% Capitulo: Integração e entrega contínuas para dados
\chapter{Integração e entrega contínuas para dados}
\label{cap:ci-cd}

Versionar o código e escrever testes só resolve metade do problema se alguém precisar se
lembrar de rodar os testes e publicar as alterações à mão. Mais cedo ou mais tarde, alguém
esquece, ou publica com pressa numa sexta-feira à tarde. A integração e a entrega contínuas
automatizam o caminho do código até a produção, de forma que cada alteração seja testada e
publicada sempre do mesmo jeito. Este capítulo trata desse caminho, do monitoramento depois
que os dados estão em produção e da documentação que se gera sozinha a partir do código.

\section{Integração contínua}
\label{sec:ci}

Integração contínua é a prática de integrar as alterações de todo o time ao código principal
com frequência, cada uma delas verificada automaticamente. Em dados, isso significa que,
quando alguém abre um pedido de revisão, um processo automático constrói os modelos alterados
num ambiente isolado, roda os testes e mostra o resultado no próprio pedido. O revisor passa a
ver não só o código, mas também se ele funciona.

Alguns cuidados tornam essa verificação útil em vez de lenta. O primeiro é construir só o que
mudou e o que depende disso, em vez do projeto inteiro, o que reduz o tempo e o custo de cada
verificação. O segundo é usar um ambiente isolado por pedido de revisão, de forma que duas
alterações em paralelo não interfiram uma na outra (\veja{sec:ambientes}). O terceiro é
incluir na verificação, além dos testes de dados, as regras de estilo e de boas práticas do
time: nomes, formatação, proibições como selecionar todas as colunas de uma tabela em modelos
de consumo. Regras que dependem da boa vontade de cada revisor são esquecidas; regras
verificadas automaticamente, não.

Com alterações escritas cada vez mais com ajuda de IA, a integração contínua ganha um papel de
rede de segurança. Ela não substitui a revisão humana, mas garante que nenhuma alteração chegue
à revisão sem ter passado pelos testes, o que libera o revisor para olhar o que só uma pessoa
consegue avaliar: se a regra está certa para o negócio.

\section{Entrega contínua}
\label{sec:cd}

Entrega contínua é o passo seguinte: depois de revisada e integrada, a alteração vai para a
produção sem nenhum passo manual \parencite{humble2010}. O mesmo processo que construiu e
testou no ambiente de integração constrói e testa em produção, na mesma ordem e com as mesmas
verificações.

A principal vantagem não é a velocidade, e sim a segurança. Publicar à mão é uma sequência de
passos que alguém precisa lembrar, e cada passo esquecido é um risco. Publicar
automaticamente é sempre igual. E, quando publicar fica fácil, as alterações ficam menores e
mais frequentes, o que torna cada uma mais fácil de revisar e, se der errado, mais fácil de
desfazer. Uma grande publicação mensal com trinta alterações é muito mais arriscada do que
trinta publicações pequenas ao longo do mês.

A pesquisa sobre desempenho de times de software reforça essa intuição. Um estudo de vários
anos com milhares de profissionais identificou quatro indicadores que distinguem os times de
melhor desempenho: a frequência com que publicam, o tempo entre uma alteração pronta e a
produção, a proporção de publicações que causam falha e o tempo para se recuperar de uma
falha. Os melhores times eram, ao mesmo tempo, mais rápidos e mais estáveis
\parencite{forsgren2018}. Os quatro indicadores se aplicam a times de dados com poucas
adaptações, e são um bom ponto de partida para medir a maturidade do processo.

\section{Observabilidade}
\label{sec:observabilidade}

Testes verificam regras conhecidas no momento da carga. Observabilidade é a capacidade de
perceber, em produção, quando algo está diferente do normal, inclusive problemas que ninguém
pensou em testar. A ideia vem do monitoramento de sistemas de software e foi adaptada para
dados, com atenção a algumas dimensões: se o dado foi atualizado no prazo, se o volume está
dentro do esperado, se a estrutura mudou, se a distribuição dos valores mudou e de onde cada
dado veio \parencite{moses2022}.

Um conceito emprestado da engenharia de confiabilidade ajuda a transformar isso em
compromisso: o objetivo de nível de serviço \parencite{beyer2016}. Em vez de prometer que os
dados estarão sempre atualizados, o time combina, por exemplo, que o painel de vendas terá
os dados do dia anterior disponíveis até as 8 horas em pelo menos 99\% dos dias úteis. O
compromisso fica mensurável, e o monitoramento passa a avisar quando ele está em risco. O
capítulo~\ref{cap:frequencia} volta a essa ideia ao discutir com que frequência os dados
precisam atualizar.

O alerta precisa chegar a quem pode agir, e na hora certa. Um alerta que dispara para o time
inteiro a cada pequena variação vira ruído e passa a ser ignorado; um que só aparece num
painel que ninguém abre não serve para nada. O objetivo é que o time descubra o problema antes
do usuário, e que o usuário, quando o problema for inevitável, seja avisado pelo time, e não
descubra sozinho.

\section{Documentação viva}
\label{sec:documentacao}

Documentação escrita à parte do código envelhece no dia em que é publicada. A alternativa é
gerar a documentação a partir do próprio código: as descrições de tabelas e colunas ficam nos
mesmos arquivos que os modelos, os testes declarados viram parte da documentação, e a
linhagem, o mapa de quais tabelas alimentam quais, é desenhada automaticamente a partir das
dependências. Quando o código muda, a documentação muda junto, porque é gerada na mesma
publicação.

Essa documentação tem hoje três leitores. O analista que precisa entender um modelo antes de
usá-lo. O revisor que precisa entender o impacto de uma alteração. E, cada vez mais, o
assistente de IA que responde perguntas sobre os dados (\veja{sec:metadados-contexto}). Para
os três, uma documentação viva, versionada e gerada a partir do código é muito mais confiável
do que um documento à parte.

\antesdepois{Publicar era um evento: alguém rodava os scripts na ordem certa, conferia à mão e
avisava o time. A documentação ficava num documento à parte, quando ficava.}
{Publicar é rotina: cada alteração revisada passa pelos testes e vai para produção sozinha. A
documentação é gerada junto e serve de contexto também para a IA.}

\rascunho{
  Relato: a implantação de um processo de integração e entrega contínuas num time de dados
  (o que foi automatizado primeiro, que problema deixou de acontecer, como a frequência de
  publicação mudou).
}

\begin{leituras}
  \leitura{humble2010}{O livro que definiu a entrega contínua. Os princípios valem para
    pipelines de dados tanto quanto para software.}
  \leitura{forsgren2018}{A pesquisa que mostrou, com dados, que velocidade e estabilidade
    andam juntas, e os quatro indicadores para medir as duas.}
  \leitura{beyer2016}{Para o conceito de objetivo de nível de serviço e a cultura de
    confiabilidade que vem com ele.}
\end{leituras}
